% ==============================================================================
% ПАРТИЯ B03-035: DISS-005-B1899 .. DISS-005-B1957
% Глава 3: Определение номиналов фазных величин и экспорт в CSV
% ==============================================================================

% DISS-005-B1899

% DISS-005-B1900
% DISS-005-TAB-084
\paragraph{Таблица 4. Алгоритм автоматического определения базисного значения фазного тока}\label{tab:ch03_phase_current_basis}

\mbox{}\par\vspace{0.2cm}
{\small
\centering
\renewcommand{\arraystretch}{1.3}
\noindent
\begin{tabular}{|>{\centering\arraybackslash}p{105pt}|>{\centering\arraybackslash}p{95pt}|>{\centering\arraybackslash}p{95pt}|>{\centering\arraybackslash}p{115pt}|}
\hline
% DISS-005-B1901
\multirow{2}{*}{Условия} &
% DISS-005-B1902
\multicolumn{2}{c|}{$H_{1(s)} > 1{,}5\,H_{x\text{-max}(s)}$} &
% DISS-005-B1903
\multirow{2}{*}{$H_{1(s)} \le 1{,}5\,H_{x\text{-max}(s)}$} \\ \cline{2-3}
% DISS-005-B1904
&
% DISS-005-B1905
$H_{1(p)} > 20$ &
% DISS-005-B1906
$H_{1(p)} \le 20$ &
% DISS-005-B1907
\\ \hline
% DISS-005-B1908
$H_{1(s)} > 30$ &
% DISS-005-B1909
\multicolumn{3}{c|}{?3} \\ \hline
% DISS-005-B1910
$0{,}03 < H_{1(s)} \le 30$ &
% DISS-005-B1911
\multicolumn{2}{c|}{$I_{\text{base}} = 5$} &
% DISS-005-B1912
?2 \\ \hline
% DISS-005-B1913
$H_{1(s)} \le 0{,}03$ &
% DISS-005-B1914
?1 &
% DISS-005-B1915
\multicolumn{2}{c|}{Noise} \\ \hline
\end{tabular}
\par}
\vspace{0.3cm}

% DISS-005-B1916

% DISS-005-B1917
% DISS-005-TAB-085
\paragraph{Таблица 5. Алгоритм автоматического определения базисного значения тока нулевой \commentnote{DISS-005-C0126}{Алексей Евдаков [2]}{Изучить код. Есть мысль – что обработка этой логики идёт не верно, и поэтому значения пропускались.}{последовательности}}\label{tab:ch03_zero_seq_current_basis}

\mbox{}\par\vspace{0.2cm}
{\small
\centering
\renewcommand{\arraystretch}{1.3}
\noindent
\begin{tabular}{|>{\centering\arraybackslash}p{340pt}|>{\centering\arraybackslash}p{80pt}|}
\hline
% DISS-005-B1918
Условия &
% DISS-005-B1919
\\ \hline
% DISS-005-B1920
$H_{1(s)} > 0{,}02$ \textbf{или} $H_{x\text{-max}(s)} > 0{,}02$ &
% DISS-005-B1921
?3 \\ \hline
% DISS-005-B1922
$0{,}02 < H_{1(s)} \le 5$ \textbf{и} $H_{x\text{-max}(s)} \le 5$ &
% DISS-005-B1923
\multirow{2}{*}{$IN_{\text{base}} = 1$} \\ \cline{1-1}
% DISS-005-B1924
$0{,}02 < H_{x\text{-max}(s)} \le 5$ \textbf{и} $H_{1(s)} \le 5$ &
% DISS-005-B1925
\\ \hline
% DISS-005-B1926
$H_{1(s)} \le 0{,}02$ \textbf{и} $H_{x\text{-max}(s)} \le 0{,}02$ &
% DISS-005-B1927
Noise \\ \hline
\end{tabular}
\par}
\vspace{0.3cm}

% DISS-005-B1928

% DISS-005-B1929
Алгоритм работает следующим образом: если H1 значительно превышает Hx\_max и находится в ожидаемом диапазоне для вторичного значения (например, около 100~В для напряжения или 5~А для фазного тока тока), то сигнал считается «чистым» и ему присваивается соответствующий стандартный вторичный базис. Если H1 мало или сопоставимо с Hx\_max, сигнал может быть классифицирован как «шум». Для тока нулевой последовательности в виду физике его возникновения алгоритм отличается. В иных случаях он определяется как спорный и требующий уточнения. Так было выявлено множество ошибок в наименованиях сигналов или осциллограммы без вторичных значений, которым были заданы за номинал~-- значения первичной сети.

% DISS-005-B1930
Затем каждому сигналу в каждой осциллограмме присваивается статус нормализации, который определяет дальнейший способ обработки. Возможные статусы:

\begin{itemize}
% DISS-005-B1931
\item \textbf{s (secondary):} Сигнал корректный, нормализация будет выполнена относительно стандартного вторичного значения (например, 100~В, 5~А). Базисное значение (\_base) будет числовым.
% DISS-005-B1932
\item \textbf{p (primary):}
% DISS-005-B1933
\item Нормализация необходимо выполнять относительно первичного значения. Этот способ используется реже, в основном, когда вторичные значения не стандартны, но перевод в них всё равно осуществляется Базисное значение (\_base) будет числовым.

\mbox{}\par
% DISS-005-B1934
\textit{\textbf{Примечание}: По возможности устанавливалась нормализация от вторичных значений, даже если это номинал в 35~кВ, это сделано для упрощений последующей нормализации}

% DISS-005-B1935
\item \textbf{Noise:} Сигнал классифицирован как шум, базисное значение (\_base) будет строкой «Noise». Как правило заменялось на схожее значение номинала на других секциях, так как сигнал точно является шумом.
% DISS-005-B1936
\item \textbf{?x:} Неопределенные статусы, потенциальные проблемы, указывающие на различные проблемы при определении базиса. Базисное значение (\_base) будет содержать соответствующий код неопределенности.
  \begin{itemize}
  % DISS-005-B1937
  \item \textbf{?1} --- возможно используются не традиционные датчики;
  % DISS-005-B1938
  \item \textbf{?2} --- возможно, из-за переходных процессов или неисправностей измерительных цепей уровень высших гармоник сопоставим или превышает уровень первой;
  % DISS-005-B1939
  \item \textbf{?3} --- как правило, указывает на отсутствие явного различия между первичными и вторичными величинами в данных файла осциллограммы (например, коэффициенты трансформации равны 1 или не заданы), либо может быть характерен для режимов с очень большими токами (сильные короткие замыкания), когда вторичные значения могут быть близки к первичным при масштабировании.
  \end{itemize}
\end{itemize}

% DISS-005-B1940
\subsubsection{Формирование таблицы коэффициентов нормализации}\label{sec:ch3_normalization_coefficients_table}

% DISS-005-B1941
Результаты описанного анализа для каждой осциллограммы и каждого унифицированного в ней сигнала сохраняются в специальную CSV-таблицу (norm\_coef.csv). Эта таблица является ключевым файлом для последующего этапа нормализации. Ее основная структура включает следующие столбцы:

\begin{itemize}
% DISS-005-B1942
\item \textbf{name:} Имя осциллограммы (MD5-хеш .dat файла).
% DISS-005-B1943
\item \textbf{norm:} Глобальный флаг агрегирует статусы нормализации отдельных сигналов осциллограммы. Он устанавливался в «NO», если, не удалось определить базис хотя бы для одного из ключевых сигналов (он получил статус неопределенности~«?»). Значение «hz\commentnote{DISS-005-C0127}{Алексей Евдаков [2]}{Ну можно поменять на Noise. Но так ведь прикольно)) Уж так сложилось))}{»} присваивается, если хотя бы один из сигналов был классифицирован как шум~-- «Noise», но осциллограмма не содержала проблем определения базиса («?») и могла быть использована.
% DISS-005-B1944
\item Далее следуют группы столбцов для каждого стандартизированного аналогового сигнала (например, для каждой из N секций шин и для каждого типа измерения Ub,~Uc,~I,~Iz), пример расписан для напряжения с секций шин (Ub):
  \begin{itemize}
  % DISS-005-B1945
  \item \textbf{\{i\}Ub\_PS:} Статус нормализации (s, p, ?1, ?2, ?3, Noise).
  % DISS-005-B1946
  \item \textbf{\{i\}Ub\_base:} Рассчитанное базисное значение (число или статус).
  % DISS-005-B1947
  \item \textbf{\{i\}Ub\_h1:} Значение H1 (амплитуда первой гармоники для оценки проведённого анализа).
  % DISS-005-B1948
  \item \textbf{\{i\}Ub\_hx:} Значение Hx\_max (максимальная амплитуда высших гармоник для оценки проведённого анализа).
  \end{itemize}
% DISS-005-B1949
\item Аналогичные группы столбцов для дифференциальных токов (\textbf{dId\_\allowbreak PS, dId\_\allowbreak base, dId\_\allowbreak h1}).
\end{itemize}

% DISS-005-B1950
Эта таблица является результатом значительной вычислительной работы и ручной коррекции (особенно для сложных и неоднозначных случаев). Наличие множества осциллограмм с некорректными или нестандартными именами сигналов, а также физическое разнообразие конфигураций энергообъектов привели к тому, что процесс создания и верификации (создание norm\_coef.csv) был итеративным и трудоемким. Было выявлено множество ошибок с именованиями и структурой исходных «.cfg» файлов.

% DISS-005-B1951
\subsection[Применение нормализации и преобразование осциллограмм в формат CSV]{Применение нормализации\\ и~преобразование осциллограмм в~формат CSV}\label{sec:ch3_csv_transformation}

% DISS-005-B1952
Последующий процесс нормализации (класс NormOsc) использует ранее созданную таблицу (norm\_coef.csv) для приведения амплитуд аналоговых сигналов в осциллограммах к относительным единицам. По уникальному имени находится соответствующий файл и если в его статусе (norm) присутствует «YES»~-- то использует нужные значения (а это по итогу осуществлено для всего датасета). Затем мгновенные значения каждого аналогового сигнала из осциллограммы делятся на его индивидуальное базисное значение (из norm\_coef.csv), предварительно умноженное на коэффициент запаса: 3 для напряжений и 20 для токов. Такое масштабирование «базиса» (фактическое увеличение делителя) осуществляется для того, чтобы результирующие нормализованные значения в номинальном режиме были значительно меньше единицы, а превышение единицы четко сигнализировало о существенных отклонениях или аппаратных аномалиях в измерительных цепях."

% DISS-005-B1953
Оговоримся, что для последующего использования в нейронных сетях, рекомендуется привести значения к диапазону от минус одного до плюс одного с нормальным распределением. Здесь функция осуществляет лишь приведение значений к единой системе координат, и поэтому для последующего использования будет требоваться «повторная» нормализация данных, которая стандартна для алгоритмов машинного обучения.

% DISS-005-B1954
Финальным этапом подготовки данных для многих задач машинного обучения является их преобразование в удобный табличный формат, такой как CSV (Comma-Separated Values). Этот процесс реализован в классе \commentnote{DISS-005-C0128}{Алексей Евдаков [2]}{Надо будет потом скорректировать по итоговой версии созданного датастеа.}{\commentnote{DISS-005-C0129}{Алексей Евдаков [2]}{Пока вообще убрал описание структуры файла}{RawToCSV}}. Используется именно CSV файлы, так как они легко читаются большинством программных средств и библиотек для анализа данных и машинного обучения (например, Pandas в Python). Табличное представление (где строки~-- это временные отсчеты, а столбцы~-- отдельные сигналы) интуитивно понятно и удобно для многих алгоритмов.

% DISS-005-B1955
В CSV файлы включаются только те аналоговые сигналы, которые были успешно унифицированы на предыдущих этапах и для которых были определены (или могли быть определены) коэффициенты нормализации. Это означает, что сигналы, не вошедшие в стандартный набор «универсальных» сигналов (раздел~3.4.4.2), отбрасываются для уменьшения «шума» в данных. Дискретные сигналы также могут быть включены.

% DISS-005-B1956
После завершения всех этих этапов мы получаем датасет осциллограмм в формате CSV, где данные стандартизированы, анонимизированы и нормализованы, что делает их пригодными для широкого спектра исследований и разработки алгоритмов в области релейной защиты и автоматики.

% DISS-005-B1957
